\documentclass[../../main.tex]{subfiles}
\graphicspath{{../../image/}{../../figure/}}
\begin{document}

\section{固定基函数的局限性}\label{section:ch06-固定基函数的局限性}

近年来，神经网络已迅速成为实际应用中最重要的机器学习技术。本书前几章已经为这一主题奠定了必要的基础：基于固定非线性基函数线性组合的回归模型可以被视为单层神经网络（参见第~4~章），同样地，基于基函数线性组合的分类模型也可以视为单层网络（参见第~5~章）。这使我们能够在讨论更复杂的多层（神经）网络之前先引入几个重要概念。

用于分类的线性基函数模型是基函数 $\{\phi_j(\mathbf{x})\}$ 的线性组合（参见第~5~章），形式为
\begin{align}
y(\mathbf{x},\mathbf{w}) = f\!\left(\sum_{j=1}^{M} w_j \phi_j(\mathbf{x}) + b\right),
\label{align:ch06-linear-basis}
\end{align}
其中 $f(\cdot)$ 是一个非线性输出激活函数；用于回归的线性模型具有相同形式，但将 $f(\cdot)$ 换成了恒等函数。由于这些基函数的一般性，这样的模型原则上可以解决任何回归或分类问题——只要其中某个基函数恰好对应所期望的输入--输出映射，可学习的线性层只需把该基函数的值复制到输出即可。

更一般地，我们期望有一个足够大且丰富的基函数集合，从而能以任意精度拟合所需函数。然而，遗憾的是，线性模型存在一些重大缺陷：它们假设基函数 $\phi_j(\mathbf{x})$ 是固定的、与训练数据无关。为了理解这些局限，我们首先观察当输入变量数量增加时线性模型的行为会如何变化。

\subsection{维度诅咒}\label{section:ch06-维度诅咒}

考虑由如下 $M$ 阶多项式给出的单输入变量简单回归模型（参见~1.2~节）：
\begin{align}
v(x,\mathbf{w}) = w_0 + w_1 x + w_2 x^2 + \cdots + w_M x^M.
\label{align:ch06-poly-1d}
\end{align}

如果增加输入变量的数量，假设有 $D$ 个输入变量 $\{x_1,\cdots,x_D\}$，那么至多三阶系数的通用多项式形式如下：
\begin{align}
v(\mathbf{x},\mathbf{w}) = w_0 + \sum_i w_i x_i + \sum_{i\le j} w_{ij} x_i x_j + \sum_{i\le j\le k} w_{ijk} x_i x_j x_k.
\label{align:ch06-poly-3d}
\end{align}

随着 $D$ 增大，独立系数的数量增长为 $O(D^3)$，而对于 $M$ 阶多项式，系数数量增长为 $O(D^M)$~\citep{Bishop2006}。我们看到，在高维空间中多项式会迅速变得笨重且几乎不具实用性。

这种朴素方法在扩展到具有更多输入变量的分类问题时同样存在许多问题。以鸢尾花数据为例：将输入空间划分为规则单元，新测试点的类别被推断为同一单元内训练点最多的类别。如果将输入空间划分为规则单元，那么单元数量将随维度 $D$ 呈指数级增长，这就要求指数级数量的训练数据才能保证单元不空，否则许多单元将不包含任何训练点，其中的测试点将得不到分类。这种在高维空间中出现的巨大困难称为\textbf{维度诅咒}（curse of dimensionality）~\citep{Bellman1961}。

\subsection{高维空间}\label{section:ch06-高维空间}

在高维空间中，我们在三维空间中形成的几何直觉可能会严重失效。考虑 $D$ 维空间中一个半径 $r=1$ 的超球体，我们想知道介于半径 $r=1-\varepsilon$ 和 $r=1$ 之间的壳层体积占超球体总体积的比例。半径为 $r$ 的 $D$ 维超球体体积 $V_D(r)$ 必须按 $r^D$ 缩放，即
\begin{align}
V_D(r) = K_D r^D,
\label{align:ch06-sphere-vol}
\end{align}
其中常数 $K_D$ 仅依赖于 $D$。该壳层占总体积的比例由下式给出（见习题~6.1）：
\begin{align}
1 - (1-\varepsilon)^D.
\label{align:ch06-shell-frac}
\end{align}

对于较大的 $D$，即使 $\varepsilon$ 很小，占比也趋近于 $1$。于是得到一个重要结论：在高维空间中，超球体的大部分体积集中在靠近表面的薄壳之中。

接下来考虑高维空间中高斯分布的行为。从笛卡尔坐标变换到极坐标并积分掉方向变量，可以得到径向密度 $p(r)$ 的表达式（见习题~6.3）。$p(r)\,\delta r$ 是位于半径 $r$ 处厚度为 $\delta r$ 的薄壳内的概率质量。对于较大的 $D$，高斯分布的大部分概率质量集中在特定半径的薄壳内。

\subsection{数据流形}\label{section:ch06-数据流形}

虽然维度诅咒确实带来了严重问题，但这并不妨碍我们找到适用高维空间的有效技术。背后的一个原因是：真实数据通常被限制在具有较低有效维度的数据空间区域内。

考虑手写数字图像。每幅图像都是高维空间中的一个点，维度由像素数决定；但由于对象可在图像内不同位置、以不同朝向出现，图像之间只有三个可变自由度。一组图像大致处在高维空间中的一个三维\textbf{流形}（manifold）上，且由于对象位置/朝向与像素强度之间的复杂关系，该流形是高度非线性的。

实际上，像素数量只是图像生成过程的人为产物。以更高分辨率捕获相同图像会增加数据空间维度 $D$，但并不改变图像仍然存在于三维流形上这一事实。如果我们能将局部化的基函数与数据流形关联起来，而不是与整个高维数据空间相关联，那么所需基函数数量将随流形维度（而非数据空间维度）呈指数级增长——而流形维度通常低得多，这将是巨大的改进。

另一种观察方式是考虑随机图像生成任务：逐像素独立从均匀分布采样 RGB 强度得到的合成图像，没有任何一幅看起来像自然图像，原因在于合成图像缺乏自然图像所表现出的相邻像素强相关性。自然图像只覆盖了整个高维空间的一小部分。

\subsection{数据依赖的基函数}\label{section:ch06-数据依赖的基函数}

独立于所解决问题而选择的简单基函数在高维空间中具有严重局限。传统方法是由专家针对每个应用手工设计特定基函数（通常称为特征），这种方法在过去数十年中是机器学习的主流。

一种数据依赖的做法是让每个训练数据点都与一个基函数相关联，从而确保基函数自动适应于底层数据流形。例如\textbf{径向基函数}（radial basis function）~\citep{BroomheadLowe1988}：每个基函数仅依赖于到中心向量的径向距离（通常是欧氏距离）。若基中心选为输入数据值 $\{\mathbf{x}_n\}$，则对每个数据点都有一个基函数，其典型形式为
\begin{align}
\phi_j(\mathbf{x}) = \exp\!\left(-\frac{\|\mathbf{x}-\mathbf{x}_j\|^2}{2s^2}\right),
\label{align:ch06-rbf}
\end{align}
其中 $s$ 控制基函数宽度。这类模型的主要问题是对大型数据集计算困难。

相关方法是\textbf{支持向量机}（Support Vector Machine, SVM）~\citep{Vapnik1995,ScholkopfSmola2002,Bishop2006}，它通过以每个训练点为中心重新定义基函数，并在训练过程中自动选择这些基函数的一个子集。最终模型中有效基函数数量通常远少于训练点数，但仍相对较大，并随训练集增大而增大；SVM 不产生概率化输出，无法自然推广到多类情况。这些方法已被深度神经网络所取代。

为了将形如~\eqref{align:ch06-linear-basis}~的线性模型应用于大规模、高维问题，我们需要一组针对所解决问题的基函数。神经网络背后的关键思想是：选择具有可学习参数的基函数 $\phi_j(\mathbf{x})$，然后在训练过程中同时调整这些参数以及系数 $\{w_j\}$。整个模型可通过基于梯度的优化方法最小化误差函数来训练（参见第~7~章），例如随机梯度下降。

有许多方法可以构造参数化的非线性基函数，但有一个关键要求：这些函数必须关于可学习参数可微，以便应用基于梯度的优化。最成功的选择是使用与~\eqref{align:ch06-linear-basis}~形式相同的基函数，其中每个基函数本身是输入的线性组合的非线性函数，而组合系数是可学习参数。这种构造显然可以递归扩展为具有多层的层次化模型，这正是深度神经网络的基础（参见~6.3~节）。

\end{document}
